\documentclass[12pt]{article}

\usepackage{amsmath,amssymb,booktabs,longtable,array,hyperref}
\topmargin -.5cm
\textheight 21cm
\oddsidemargin -.125cm
\textwidth 17cm

\newcommand{\SO}{\mathrm{SO}}
\newcommand{\SU}{\mathrm{SU}}
\newcommand{\D}{\Delta}
\newcommand{\la}{\langle}
\newcommand{\ra}{\rangle}
\newcommand{\bb}{\mathbf}
\newcommand{\half}{\frac{1}{2}}

\begin{document}
\sloppy

\begin{center}
{\Large\bf Mass-Level-1 Primary States and the Konishi Table}
\end{center}

\section{Goal}

The purpose of this note is to compare, at the level of \emph{primary-state labels},
the full first-massive flat-space Type IIB spectrum with the states listed in Table~1 of
\texttt{references/1102.1209/1102.1209.tex}.

The comparison is performed with the following rules.
\begin{itemize}
\item On the AdS side, I keep the primary labels $(s_L,s_R)$ exactly as they are shown in Table~1.
\item On the sphere side, I branch every $\SO(6)$ irrep in Table~1 to $\SO(5)$.
\item On the flat-space side, I decompose the first-massive physical representations under
$\SO(1,4)\times \SO(5)$ and then keep the corresponding compact $\SO(4)$ spin labels
$(s_L,s_R)$.
\item The comparison key is therefore
\begin{equation}
\big((s_L,s_R),\,\la a,b\ra\big),
\end{equation}
with multiplicities retained.
\end{itemize}

The independent computation is implemented in
\texttt{scripts/primary\_branching\_scan.py}. Existing notes were used only \emph{after}
the derivation as a validation check.

\section{Independent Branching Data}

\subsection{Left-moving first-massive building blocks}

Write the vector of the little group as
\begin{equation}
\bb 9 \to \Big(\half,\half;\la 0,0\ra\Big)\oplus \Big(0,0;\la 0,1\ra\Big),
\end{equation}
where $\la a,b\ra$ denotes a $\SO(5)$ irrep in $B_2$ Dynkin notation.

The first-massive physical left-moving representations are
\begin{equation}
\bb{44}=\mathrm{Sym}^2_0(\bb 9),\qquad
\bb{84}=\Lambda^3(\bb 9),\qquad
\bb{128}\subset \bb 9\otimes \bb{16}.
\end{equation}
Their independent decomposition is:

\begin{table}[h]
\centering
\begin{tabular}{>{$\displaystyle}c<{$} >{$\displaystyle}l<{$}}
\toprule
\text{rep} & \SO(4)\times \SO(5)\text{ decomposition} \\
\midrule
\bb{44} &
(0,0;\la 0,0\ra)\oplus (0,0;\la 0,2\ra)\oplus (\half,\half;\la 0,1\ra)\oplus (1,1;\la 0,0\ra)
\\[2mm]
\bb{84} &
(\half,\half;\la 0,0\ra)\oplus (1,0;\la 0,1\ra)\oplus (0,1;\la 0,1\ra)\oplus
(\half,\half;\la 2,0\ra)\oplus (0,0;\la 2,0\ra)
\\[2mm]
\bb{128} &
(1,\half;\la 1,0\ra)\oplus (\half,1;\la 1,0\ra)\oplus (\half,0;\la 1,0\ra)\oplus
(0,\half;\la 1,0\ra)
\\ &
\oplus (\half,0;\la 1,1\ra)\oplus (0,\half;\la 1,1\ra)
\\
\bottomrule
\end{tabular}
\caption{Independent decomposition of the left-moving first-massive building blocks.}
\end{table}

The closed-string sectors are then formed as
\begin{equation}
\text{NSNS}=(\bb{44}\oplus \bb{84})_L\otimes (\bb{44}\oplus \bb{84})_R,\qquad
\text{RR}=\bb{128}_L\otimes \bb{128}_R,
\end{equation}
and similarly for the mixed sectors NSR and RNS.

\subsection{Overall multiplicity counts}

After tensoring left and right movers and decomposing all $\SO(5)$ products independently,
the total primary-state multiplicity counts are:

\begin{table}[h]
\centering
\begin{tabular}{lc}
\toprule
\text{sector} & \text{total multiplicity} \\
\midrule
\text{NSNS} & 306 \\
\text{RR} & 298 \\
\text{NSR} & 242 \\
\text{RNS} & 242 \\
\midrule
\text{Total flat-space multiplicity} & 1088 \\
\bottomrule
\end{tabular}
\caption{Multiplicity counts for the full first-massive flat-space primary spectrum.}
\end{table}

\section{Comparison with the Konishi Table}

Every $\SO(6)$ irrep in Table~1 was independently branched to $\SO(5)$, while the displayed
primary spin labels $(s_L,s_R)$ were kept unchanged.

The global comparison is:

\begin{table}[h]
\centering
\begin{tabular}{lc}
\toprule
\text{quantity} & \text{value} \\
\midrule
\text{Distinct flat-space irreps} & 109 \\
\text{Distinct branched Konishi irreps} & 109 \\
\text{Set-level intersection} & 109 \\
\text{Flat-only distinct irreps} & 2 \\
\text{Konishi-only distinct irreps} & 0 \\
\midrule
\text{Total flat-space multiplicity} & 1088 \\
\text{Total branched Konishi multiplicity} & 1086 \\
\bottomrule
\end{tabular}
\caption{Global comparison of primary-state labels $\big((s_L,s_R),\la a,b\ra\big)$.}
\end{table}

So the answer is:
\begin{itemize}
\item At the level of \emph{distinct irreps}, every flat-space mass-level-1 primary irrep appears in the
branched Konishi table.
\item The only failure is at the level of multiplicities.
\end{itemize}

The complete mismatch is:

\begin{table}[h]
\centering
\begin{tabular}{ccc}
\toprule
\text{irrep} & \text{flat multiplicity} & \text{Konishi multiplicity} \\
\midrule
$(0,\half;\la 1,0\ra)$ & 32 & 31 \\
$(\half,0;\la 1,0\ra)$ & 32 & 31 \\
\bottomrule
\end{tabular}
\caption{The only multiplicity mismatch in the comparison.}
\end{table}

There are no irreps that appear on the Konishi side but fail to appear on the flat-space side.

\subsection{Explanation of the result}

The comparison is surprisingly tight:
\begin{equation}
\text{all 109 distinct flat-space irreps occur in the branched Konishi table.}
\end{equation}
Thus the branched Konishi table is large enough to cover the entire first-massive primary-state
content at the level of representation \emph{types}.

The only discrepancy is that the two chiral spinor-vector sphere-spinor labels
\begin{equation}
(0,\half;\la 1,0\ra),\qquad (\half,0;\la 1,0\ra)
\end{equation}
appear once fewer on the Konishi side than in the full flat-space count. In other words,
the mismatch is not structural. It does not remove an entire representation family.
It is a very small multiplicity deficit in precisely two fermionic channels.

For the present primary-state question, the clean summary is therefore:
\begin{equation}
\text{set-level match: yes, \qquad multiplicity-level match: no, but only by }2.
\end{equation}

\subsection{Source of the mismatch}

The mismatch can be localized very sharply.
\begin{itemize}
\item It is \emph{not} caused by the $\SO(6)\to \SO(5)$ branching. That part agrees exactly with the
independent check against the existing note.
\item It is \emph{not} caused by the decomposition of the $\bb{44}$, $\bb{84}$, or $\bb{128}$ under
$\SO(4)\times \SO(5)$. Those also agree with the existing note.
\item The two excess copies on the flat-space side live entirely in the mixed fermionic sectors NSR and RNS.
\end{itemize}

More concretely, each of the problematic irreps receives
\begin{equation}
16\ \text{copies from NSR} \qquad\text{and}\qquad 16\ \text{copies from RNS},
\end{equation}
so the flat-space total is $32$, whereas the branched Konishi table contains only $31$.

This points to the actual cause: on the flat-space side I counted \emph{all} compact
$\SO(4)$ pieces obtained from the first-massive mixed sectors, while on the Konishi side
Table~1 counts only superconformal primaries. For the two irreps
\begin{equation}
(0,\half;\la 1,0\ra),\qquad (\half,0;\la 1,0\ra),
\end{equation}
the naive flat-space compact reduction appears to overcount by one linear combination in each
chirality. The present script does \emph{not} implement the full AdS primary projection, so this
should be read as an interpretation of the data rather than as a separate theorem.

What \emph{is} established directly by the computation is that the phenomenon is isolated.
Among all mixed-sector irreps with sphere label $\la 1,0\ra$, the complete tower is
\begin{equation}
\begin{aligned}
&(0,\half),(0,\tfrac32),(\half,0),(\half,1),(\half,2),(1,\half),\\
&(1,\tfrac32),(\tfrac32,0),(\tfrac32,1),(\tfrac32,2),(2,\half),(2,\tfrac32),
\end{aligned}
\end{equation}
and every entry in this tower matches the Konishi multiplicity exactly except the two bottom
states $(0,\half;\la 1,0\ra)$ and $(\half,0;\la 1,0\ra)$.
They are also the only mixed $\la 1,0\ra$ states whose lowest conformal weight is
$\Delta_0=\frac52$; all higher-spin states in the same family start at $\Delta_0\ge \frac72$ and
agree perfectly with Konishi.

So this pair is special because it sits at the very bottom of the mixed fermionic
$\la 1,0\ra$ tower: minimal $\SO(4)$ spin and minimal $\Delta_0$.
That is exactly where one expects a naive compact-state count to be most sensitive to the
difference between ``all compact components'' and ``actual AdS primaries''.

Equivalently, if one corrects the flat-space primary count by subtracting one copy of
\begin{equation}
(0,\half;\la 1,0\ra)\qquad\text{and}\qquad (\half,0;\la 1,0\ra),
\end{equation}
then the multiplicities match \emph{exactly}:
\begin{equation}
1088-2=1086.
\end{equation}

So the mismatch is best understood not as a genuine disagreement in representation content,
but as the residual effect of using a coarse flat-space compact-spin projection as a proxy
for the actual AdS primary-state condition. The important point for the present comparison is that
no analogous hidden mismatch shows up anywhere else in the spectrum: every other irrep already
matches before any such correction.

\section{Check Against the Existing Note}

After the independent derivation, I compared the results against
\texttt{notes/konishi\_linearized\_ansatz.tex}. The overlap agrees exactly:

\begin{table}[h]
\centering
\begin{tabular}{lc}
\toprule
\text{check} & \text{mismatches} \\
\midrule
$\SO(6)\to \SO(5)$ branching table & 0 \\
flat $\bb{44},\bb{84},\bb{128}$ decompositions & 0 \\
\bottomrule
\end{tabular}
\caption{Validation against the existing note after the independent derivation.}
\end{table}

\appendix

\section{Independent Group-Theory Algorithm}

\subsection{Branching \texorpdfstring{$\SO(6)\to \SO(5)$}{SO(6)->SO(5)}}

I treat $\SO(6)$ as $D_3$ and $\SO(5)$ as $B_2$.
For an $\SO(6)\simeq \SU(4)$ irrep with Dynkin labels
\begin{equation}
[a,b,c],
\end{equation}
I convert to the scaled orthogonal highest weight
\begin{equation}
(\lambda_1,\lambda_2,\lambda_3)=(a+c+2b,\ a+c,\ c-a).
\end{equation}

The $D_3\to B_2$ branching is then obtained from the orthogonal interlacing rule:
\begin{equation}
\lambda_1\ge m_1\ge \lambda_2\ge m_2\ge |\lambda_3|,
\end{equation}
with $m_1,m_2$ of the same parity as the $\lambda_i$ in the present scaled convention.
Each admissible pair $(m_1,m_2)$ contributes one $\SO(5)$ irrep with Dynkin labels
\begin{equation}
\la a,b\ra=\la m_2,\ \tfrac{m_1-m_2}{2}\ra.
\end{equation}
This directly produces the full $\SO(6)\to \SO(5)$ branching table used in the note.

\subsection{Tensor products in \texorpdfstring{$\SO(5)$}{SO(5)}}

To multiply $\SO(5)$ irreps, I compute the full weight system of a highest weight
\begin{equation}
\la a,b\ra
\end{equation}
using Freudenthal recursion for the $B_2$ root system. In scaled orthogonal coordinates,
the positive roots are
\begin{equation}
(2,-2),\qquad (0,2),\qquad (2,0),\qquad (2,2),
\end{equation}
and the Weyl vector is
\begin{equation}
\rho=(3,1).
\end{equation}

Starting from the highest weight, I generate lower weights by subtracting simple roots and
determine their multiplicities recursively from the Freudenthal formula. After that, tensor
products are computed in two steps:
\begin{enumerate}
\item add left and right weights with multiplicities;
\item repeatedly subtract irreducible $B_2$ characters, starting from the dominant highest weight
still present.
\end{enumerate}
This yields the complete $\SO(5)$ tensor-product decomposition with multiplicities.

\subsection{Flat-space first-massive spectrum}

The flat-space side is constructed from the physical first-massive little-group representations.

\paragraph{The bosons.}
Decompose the vector as
\begin{equation}
\bb 9\to \Big(\half,\half;\la 0,0\ra\Big)\oplus \Big(0,0;\la 0,1\ra\Big).
\end{equation}
Then compute
\begin{equation}
\bb{44}=\mathrm{Sym}^2_0(\bb 9),\qquad \bb{84}=\Lambda^3(\bb 9)
\end{equation}
by standard tensor algebra.

\paragraph{The fermions.}
The spinor branches as
\begin{equation}
\bb{16}\to \Big(\half,0;\la 1,0\ra\Big)\oplus \Big(0,\half;\la 1,0\ra\Big).
\end{equation}
The vector-spinor is obtained from
\begin{equation}
\bb 9\otimes \bb{16}
\end{equation}
and then the gamma-trace $\bb{16}$ is removed, leaving the gamma-traceless
representation $\bb{128}$.

\paragraph{Closed-string sectors.}
Finally, I tensor left and right movers to obtain the NSNS, RR, NSR, and RNS sectors, and
collect the resulting primary-state labels
\begin{equation}
\big((s_L,s_R),\la a,b\ra\big)
\end{equation}
with multiplicity.

\subsection{Comparison stage}

Each Table-1 entry is parsed as
\begin{equation}
(\D_0,\ [p,q,r],\ (s_L,s_R),\ \text{multiplicity}),
\end{equation}
branched independently on the $\SO(6)$ side, and accumulated into the same key
\begin{equation}
\big((s_L,s_R),\la a,b\ra\big).
\end{equation}
The final output compares the flat-space and Konishi multisets directly, so both set-level
and multiplicity-level statements follow immediately from the same data.

\clearpage
\section{Complete Flat-Space Irrep Table}

For completeness, I include the full list of primary-state irreps obtained from reducing the
first-massive flat-space spectrum. The table is grouped by compact spin label $(s_L,s_R)$ and
also shows the sector split into NSNS, RR, and mixed NSR/RNS contributions.

\begingroup
\begin{longtable}{c c c c c c}
\caption{Complete flat-space primary-state irreps after reducing the first-massive spectrum, grouped by compact spin label.}\label{tab:flat-full}\\
\toprule
$(s_L,s_R)$ & $\SO(5)$ irrep & total & NSNS & RR & mixed \\
\midrule
\endfirsthead
\toprule
$(s_L,s_R)$ & $\SO(5)$ irrep & total & NSNS & RR & mixed \\
\midrule
\endhead
\multicolumn{6}{l}{\textbf{Spin $(0,0)$}}\\
$(0,0)$ & $\langle 0,0\rangle$ & 15 & 9 & 6 & 0 \\
$(0,0)$ & $\langle 0,1\rangle$ & 16 & 6 & 10 & 0 \\
$(0,0)$ & $\langle 2,0\rangle$ & 26 & 14 & 12 & 0 \\
$(0,0)$ & $\langle 0,2\rangle$ & 16 & 10 & 6 & 0 \\
$(0,0)$ & $\langle 2,1\rangle$ & 14 & 6 & 8 & 0 \\
$(0,0)$ & $\langle 4,0\rangle$ & 5 & 3 & 2 & 0 \\
$(0,0)$ & $\langle 0,3\rangle$ & 2 & 0 & 2 & 0 \\
$(0,0)$ & $\langle 2,2\rangle$ & 5 & 3 & 2 & 0 \\
$(0,0)$ & $\langle 0,4\rangle$ & 1 & 1 & 0 & 0 \\
\midrule
\multicolumn{6}{l}{\textbf{Spin $(0,1/2)$}}\\
$(0,1/2)$ & $\langle 1,0\rangle$ & 32 & 0 & 0 & 32 \\
$(0,1/2)$ & $\langle 1,1\rangle$ & 34 & 0 & 0 & 34 \\
$(0,1/2)$ & $\langle 3,0\rangle$ & 16 & 0 & 0 & 16 \\
$(0,1/2)$ & $\langle 1,2\rangle$ & 12 & 0 & 0 & 12 \\
$(0,1/2)$ & $\langle 3,1\rangle$ & 6 & 0 & 0 & 6 \\
$(0,1/2)$ & $\langle 1,3\rangle$ & 2 & 0 & 0 & 2 \\
\midrule
\multicolumn{6}{l}{\textbf{Spin $(0,1)$}}\\
$(0,1)$ & $\langle 0,0\rangle$ & 11 & 5 & 6 & 0 \\
$(0,1)$ & $\langle 0,1\rangle$ & 23 & 13 & 10 & 0 \\
$(0,1)$ & $\langle 2,0\rangle$ & 20 & 9 & 11 & 0 \\
$(0,1)$ & $\langle 0,2\rangle$ & 8 & 3 & 5 & 0 \\
$(0,1)$ & $\langle 2,1\rangle$ & 13 & 7 & 6 & 0 \\
$(0,1)$ & $\langle 4,0\rangle$ & 2 & 1 & 1 & 0 \\
$(0,1)$ & $\langle 0,3\rangle$ & 3 & 2 & 1 & 0 \\
$(0,1)$ & $\langle 2,2\rangle$ & 1 & 0 & 1 & 0 \\
\midrule
\multicolumn{6}{l}{\textbf{Spin $(0,3/2)$}}\\
$(0,3/2)$ & $\langle 1,0\rangle$ & 12 & 0 & 0 & 12 \\
$(0,3/2)$ & $\langle 1,1\rangle$ & 8 & 0 & 0 & 8 \\
$(0,3/2)$ & $\langle 3,0\rangle$ & 4 & 0 & 0 & 4 \\
$(0,3/2)$ & $\langle 1,2\rangle$ & 2 & 0 & 0 & 2 \\
\midrule
\multicolumn{6}{l}{\textbf{Spin $(0,2)$}}\\
$(0,2)$ & $\langle 0,0\rangle$ & 3 & 2 & 1 & 0 \\
$(0,2)$ & $\langle 0,1\rangle$ & 1 & 0 & 1 & 0 \\
$(0,2)$ & $\langle 2,0\rangle$ & 2 & 1 & 1 & 0 \\
$(0,2)$ & $\langle 0,2\rangle$ & 1 & 1 & 0 & 0 \\
\midrule
\multicolumn{6}{l}{\textbf{Spin $(1/2,0)$}}\\
$(1/2,0)$ & $\langle 1,0\rangle$ & 32 & 0 & 0 & 32 \\
$(1/2,0)$ & $\langle 1,1\rangle$ & 34 & 0 & 0 & 34 \\
$(1/2,0)$ & $\langle 3,0\rangle$ & 16 & 0 & 0 & 16 \\
$(1/2,0)$ & $\langle 1,2\rangle$ & 12 & 0 & 0 & 12 \\
$(1/2,0)$ & $\langle 3,1\rangle$ & 6 & 0 & 0 & 6 \\
$(1/2,0)$ & $\langle 1,3\rangle$ & 2 & 0 & 0 & 2 \\
\midrule
\multicolumn{6}{l}{\textbf{Spin $(1/2,1/2)$}}\\
$(1/2,1/2)$ & $\langle 0,0\rangle$ & 20 & 10 & 10 & 0 \\
$(1/2,1/2)$ & $\langle 0,1\rangle$ & 36 & 18 & 18 & 0 \\
$(1/2,1/2)$ & $\langle 2,0\rangle$ & 40 & 20 & 20 & 0 \\
$(1/2,1/2)$ & $\langle 0,2\rangle$ & 20 & 10 & 10 & 0 \\
$(1/2,1/2)$ & $\langle 2,1\rangle$ & 24 & 12 & 12 & 0 \\
$(1/2,1/2)$ & $\langle 4,0\rangle$ & 4 & 2 & 2 & 0 \\
$(1/2,1/2)$ & $\langle 0,3\rangle$ & 4 & 2 & 2 & 0 \\
$(1/2,1/2)$ & $\langle 2,2\rangle$ & 4 & 2 & 2 & 0 \\
\midrule
\multicolumn{6}{l}{\textbf{Spin $(1/2,1)$}}\\
$(1/2,1)$ & $\langle 1,0\rangle$ & 32 & 0 & 0 & 32 \\
$(1/2,1)$ & $\langle 1,1\rangle$ & 30 & 0 & 0 & 30 \\
$(1/2,1)$ & $\langle 3,0\rangle$ & 12 & 0 & 0 & 12 \\
$(1/2,1)$ & $\langle 1,2\rangle$ & 8 & 0 & 0 & 8 \\
$(1/2,1)$ & $\langle 3,1\rangle$ & 2 & 0 & 0 & 2 \\
\midrule
\multicolumn{6}{l}{\textbf{Spin $(1/2,3/2)$}}\\
$(1/2,3/2)$ & $\langle 0,0\rangle$ & 8 & 4 & 4 & 0 \\
$(1/2,3/2)$ & $\langle 0,1\rangle$ & 12 & 6 & 6 & 0 \\
$(1/2,3/2)$ & $\langle 2,0\rangle$ & 12 & 6 & 6 & 0 \\
$(1/2,3/2)$ & $\langle 0,2\rangle$ & 4 & 2 & 2 & 0 \\
$(1/2,3/2)$ & $\langle 2,1\rangle$ & 4 & 2 & 2 & 0 \\
\midrule
\multicolumn{6}{l}{\textbf{Spin $(1/2,2)$}}\\
$(1/2,2)$ & $\langle 1,0\rangle$ & 4 & 0 & 0 & 4 \\
$(1/2,2)$ & $\langle 1,1\rangle$ & 2 & 0 & 0 & 2 \\
\midrule
\multicolumn{6}{l}{\textbf{Spin $(1,0)$}}\\
$(1,0)$ & $\langle 0,0\rangle$ & 11 & 5 & 6 & 0 \\
$(1,0)$ & $\langle 0,1\rangle$ & 23 & 13 & 10 & 0 \\
$(1,0)$ & $\langle 2,0\rangle$ & 20 & 9 & 11 & 0 \\
$(1,0)$ & $\langle 0,2\rangle$ & 8 & 3 & 5 & 0 \\
$(1,0)$ & $\langle 2,1\rangle$ & 13 & 7 & 6 & 0 \\
$(1,0)$ & $\langle 4,0\rangle$ & 2 & 1 & 1 & 0 \\
$(1,0)$ & $\langle 0,3\rangle$ & 3 & 2 & 1 & 0 \\
$(1,0)$ & $\langle 2,2\rangle$ & 1 & 0 & 1 & 0 \\
\midrule
\multicolumn{6}{l}{\textbf{Spin $(1,1/2)$}}\\
$(1,1/2)$ & $\langle 1,0\rangle$ & 32 & 0 & 0 & 32 \\
$(1,1/2)$ & $\langle 1,1\rangle$ & 30 & 0 & 0 & 30 \\
$(1,1/2)$ & $\langle 3,0\rangle$ & 12 & 0 & 0 & 12 \\
$(1,1/2)$ & $\langle 1,2\rangle$ & 8 & 0 & 0 & 8 \\
$(1,1/2)$ & $\langle 3,1\rangle$ & 2 & 0 & 0 & 2 \\
\midrule
\multicolumn{6}{l}{\textbf{Spin $(1,1)$}}\\
$(1,1)$ & $\langle 0,0\rangle$ & 14 & 8 & 6 & 0 \\
$(1,1)$ & $\langle 0,1\rangle$ & 19 & 9 & 10 & 0 \\
$(1,1)$ & $\langle 2,0\rangle$ & 20 & 10 & 10 & 0 \\
$(1,1)$ & $\langle 0,2\rangle$ & 10 & 6 & 4 & 0 \\
$(1,1)$ & $\langle 2,1\rangle$ & 7 & 3 & 4 & 0 \\
$(1,1)$ & $\langle 4,0\rangle$ & 1 & 1 & 0 & 0 \\
\midrule
\multicolumn{6}{l}{\textbf{Spin $(1,3/2)$}}\\
$(1,3/2)$ & $\langle 1,0\rangle$ & 12 & 0 & 0 & 12 \\
$(1,3/2)$ & $\langle 1,1\rangle$ & 8 & 0 & 0 & 8 \\
$(1,3/2)$ & $\langle 3,0\rangle$ & 2 & 0 & 0 & 2 \\
\midrule
\multicolumn{6}{l}{\textbf{Spin $(1,2)$}}\\
$(1,2)$ & $\langle 0,0\rangle$ & 2 & 1 & 1 & 0 \\
$(1,2)$ & $\langle 0,1\rangle$ & 3 & 2 & 1 & 0 \\
$(1,2)$ & $\langle 2,0\rangle$ & 1 & 0 & 1 & 0 \\
\midrule
\multicolumn{6}{l}{\textbf{Spin $(3/2,0)$}}\\
$(3/2,0)$ & $\langle 1,0\rangle$ & 12 & 0 & 0 & 12 \\
$(3/2,0)$ & $\langle 1,1\rangle$ & 8 & 0 & 0 & 8 \\
$(3/2,0)$ & $\langle 3,0\rangle$ & 4 & 0 & 0 & 4 \\
$(3/2,0)$ & $\langle 1,2\rangle$ & 2 & 0 & 0 & 2 \\
\midrule
\multicolumn{6}{l}{\textbf{Spin $(3/2,1/2)$}}\\
$(3/2,1/2)$ & $\langle 0,0\rangle$ & 8 & 4 & 4 & 0 \\
$(3/2,1/2)$ & $\langle 0,1\rangle$ & 12 & 6 & 6 & 0 \\
$(3/2,1/2)$ & $\langle 2,0\rangle$ & 12 & 6 & 6 & 0 \\
$(3/2,1/2)$ & $\langle 0,2\rangle$ & 4 & 2 & 2 & 0 \\
$(3/2,1/2)$ & $\langle 2,1\rangle$ & 4 & 2 & 2 & 0 \\
\midrule
\multicolumn{6}{l}{\textbf{Spin $(3/2,1)$}}\\
$(3/2,1)$ & $\langle 1,0\rangle$ & 12 & 0 & 0 & 12 \\
$(3/2,1)$ & $\langle 1,1\rangle$ & 8 & 0 & 0 & 8 \\
$(3/2,1)$ & $\langle 3,0\rangle$ & 2 & 0 & 0 & 2 \\
\midrule
\multicolumn{6}{l}{\textbf{Spin $(3/2,3/2)$}}\\
$(3/2,3/2)$ & $\langle 0,0\rangle$ & 4 & 2 & 2 & 0 \\
$(3/2,3/2)$ & $\langle 0,1\rangle$ & 4 & 2 & 2 & 0 \\
$(3/2,3/2)$ & $\langle 2,0\rangle$ & 4 & 2 & 2 & 0 \\
\midrule
\multicolumn{6}{l}{\textbf{Spin $(3/2,2)$}}\\
$(3/2,2)$ & $\langle 1,0\rangle$ & 2 & 0 & 0 & 2 \\
\midrule
\multicolumn{6}{l}{\textbf{Spin $(2,0)$}}\\
$(2,0)$ & $\langle 0,0\rangle$ & 3 & 2 & 1 & 0 \\
$(2,0)$ & $\langle 0,1\rangle$ & 1 & 0 & 1 & 0 \\
$(2,0)$ & $\langle 2,0\rangle$ & 2 & 1 & 1 & 0 \\
$(2,0)$ & $\langle 0,2\rangle$ & 1 & 1 & 0 & 0 \\
\midrule
\multicolumn{6}{l}{\textbf{Spin $(2,1/2)$}}\\
$(2,1/2)$ & $\langle 1,0\rangle$ & 4 & 0 & 0 & 4 \\
$(2,1/2)$ & $\langle 1,1\rangle$ & 2 & 0 & 0 & 2 \\
\midrule
\multicolumn{6}{l}{\textbf{Spin $(2,1)$}}\\
$(2,1)$ & $\langle 0,0\rangle$ & 2 & 1 & 1 & 0 \\
$(2,1)$ & $\langle 0,1\rangle$ & 3 & 2 & 1 & 0 \\
$(2,1)$ & $\langle 2,0\rangle$ & 1 & 0 & 1 & 0 \\
\midrule
\multicolumn{6}{l}{\textbf{Spin $(2,3/2)$}}\\
$(2,3/2)$ & $\langle 1,0\rangle$ & 2 & 0 & 0 & 2 \\
\midrule
\multicolumn{6}{l}{\textbf{Spin $(2,2)$}}\\
$(2,2)$ & $\langle 0,0\rangle$ & 1 & 1 & 0 & 0 \\
\bottomrule
\end{longtable}

\endgroup

\clearpage
\section{Complete Konishi Irrep Table}

This appendix contains the full list of primary-state irreps obtained from the Konishi table
after branching $\SO(6)\to \SO(5)$. The table is again grouped by compact spin label and each
row is graded by the corresponding values of $\D_0$.

\begingroup
\begin{longtable}{c c c p{8.8cm}}
\caption{Complete Konishi primary-state irreps after branching $\SO(6)\to \SO(5)$, organized by compact spin label and graded by $\Delta_0$.}\label{tab:konishi-full}\\
\toprule
$(s_L,s_R)$ & $\SO(5)$ irrep & total & $\Delta_0$ grading \\
\midrule
\endfirsthead
\toprule
$(s_L,s_R)$ & $\SO(5)$ irrep & total & $\Delta_0$ grading \\
\midrule
\endhead
\multicolumn{4}{l}{\textbf{Spin $(0,0)$}}\\
$(0,0)$ & $\langle 0,0\rangle$ & 15 & $2$: 1, $4$: 3, $6$: 7, $8$: 3, $\frac{19}{2}$: 1 \\
$(0,0)$ & $\langle 0,1\rangle$ & 16 & $4$: 3, $5$: 2, $6$: 6, $7$: 2, $8$: 3 \\
$(0,0)$ & $\langle 2,0\rangle$ & 26 & $3$: 2, $4$: 1, $5$: 8, $6$: 4, $7$: 8, $8$: 1, $9$: 2 \\
$(0,0)$ & $\langle 0,2\rangle$ & 16 & $4$: 3, $5$: 2, $6$: 6, $7$: 2, $8$: 3 \\
$(0,0)$ & $\langle 2,1\rangle$ & 14 & $4$: 1, $5$: 4, $6$: 4, $7$: 4, $8$: 1 \\
$(0,0)$ & $\langle 4,0\rangle$ & 5 & $4$: 1, $6$: 3, $8$: 1 \\
$(0,0)$ & $\langle 0,3\rangle$ & 2 & $6$: 2 \\
$(0,0)$ & $\langle 2,2\rangle$ & 5 & $5$: 2, $6$: 1, $7$: 2 \\
$(0,0)$ & $\langle 0,4\rangle$ & 1 & $6$: 1 \\
\midrule
\multicolumn{4}{l}{\textbf{Spin $(0,1/2)$}}\\
$(0,1/2)$ & $\langle 1,0\rangle$ & 31 & $\frac{5}{2}$: 1, $\frac{7}{2}$: 2, $\frac{9}{2}$: 5, $\frac{11}{2}$: 8, $\frac{13}{2}$: 8, $\frac{15}{2}$: 5, $\frac{17}{2}$: 2 \\
$(0,1/2)$ & $\langle 1,1\rangle$ & 34 & $\frac{7}{2}$: 2, $\frac{9}{2}$: 6, $\frac{11}{2}$: 9, $\frac{13}{2}$: 9, $\frac{15}{2}$: 6, $\frac{17}{2}$: 2 \\
$(0,1/2)$ & $\langle 3,0\rangle$ & 16 & $\frac{7}{2}$: 1, $\frac{9}{2}$: 2, $\frac{11}{2}$: 5, $\frac{13}{2}$: 5, $\frac{15}{2}$: 2, $\frac{17}{2}$: 1 \\
$(0,1/2)$ & $\langle 1,2\rangle$ & 12 & $\frac{9}{2}$: 2, $\frac{11}{2}$: 4, $\frac{13}{2}$: 4, $\frac{15}{2}$: 2 \\
$(0,1/2)$ & $\langle 3,1\rangle$ & 6 & $\frac{9}{2}$: 1, $\frac{11}{2}$: 2, $\frac{13}{2}$: 2, $\frac{15}{2}$: 1 \\
$(0,1/2)$ & $\langle 1,3\rangle$ & 2 & $\frac{11}{2}$: 1, $\frac{13}{2}$: 1 \\
\midrule
\multicolumn{4}{l}{\textbf{Spin $(0,1)$}}\\
$(0,1)$ & $\langle 0,0\rangle$ & 11 & $3$: 1, $5$: 4, $6$: 1, $7$: 4, $9$: 1 \\
$(0,1)$ & $\langle 0,1\rangle$ & 23 & $3$: 1, $4$: 2, $5$: 6, $6$: 5, $7$: 6, $8$: 2, $9$: 1 \\
$(0,1)$ & $\langle 2,0\rangle$ & 20 & $4$: 3, $5$: 3, $6$: 8, $7$: 3, $8$: 3 \\
$(0,1)$ & $\langle 0,2\rangle$ & 8 & $5$: 3, $6$: 2, $7$: 3 \\
$(0,1)$ & $\langle 2,1\rangle$ & 13 & $4$: 1, $5$: 3, $6$: 5, $7$: 3, $8$: 1 \\
$(0,1)$ & $\langle 4,0\rangle$ & 2 & $5$: 1, $7$: 1 \\
$(0,1)$ & $\langle 0,3\rangle$ & 3 & $5$: 1, $6$: 1, $7$: 1 \\
$(0,1)$ & $\langle 2,2\rangle$ & 1 & $6$: 1 \\
\midrule
\multicolumn{4}{l}{\textbf{Spin $(0,3/2)$}}\\
$(0,3/2)$ & $\langle 1,0\rangle$ & 12 & $\frac{7}{2}$: 1, $\frac{9}{2}$: 2, $\frac{11}{2}$: 3, $\frac{13}{2}$: 3, $\frac{15}{2}$: 2, $\frac{17}{2}$: 1 \\
$(0,3/2)$ & $\langle 1,1\rangle$ & 8 & $\frac{9}{2}$: 1, $\frac{11}{2}$: 3, $\frac{13}{2}$: 3, $\frac{15}{2}$: 1 \\
$(0,3/2)$ & $\langle 3,0\rangle$ & 4 & $\frac{9}{2}$: 1, $\frac{11}{2}$: 1, $\frac{13}{2}$: 1, $\frac{15}{2}$: 1 \\
$(0,3/2)$ & $\langle 1,2\rangle$ & 2 & $\frac{11}{2}$: 1, $\frac{13}{2}$: 1 \\
\midrule
\multicolumn{4}{l}{\textbf{Spin $(0,2)$}}\\
$(0,2)$ & $\langle 0,0\rangle$ & 3 & $4$: 1, $6$: 1, $8$: 1 \\
$(0,2)$ & $\langle 0,1\rangle$ & 1 & $6$: 1 \\
$(0,2)$ & $\langle 2,0\rangle$ & 2 & $5$: 1, $7$: 1 \\
$(0,2)$ & $\langle 0,2\rangle$ & 1 & $6$: 1 \\
\midrule
\multicolumn{4}{l}{\textbf{Spin $(1/2,0)$}}\\
$(1/2,0)$ & $\langle 1,0\rangle$ & 31 & $\frac{5}{2}$: 1, $\frac{7}{2}$: 2, $\frac{9}{2}$: 5, $\frac{11}{2}$: 8, $\frac{13}{2}$: 8, $\frac{15}{2}$: 5, $\frac{17}{2}$: 2 \\
$(1/2,0)$ & $\langle 1,1\rangle$ & 34 & $\frac{7}{2}$: 2, $\frac{9}{2}$: 6, $\frac{11}{2}$: 9, $\frac{13}{2}$: 9, $\frac{15}{2}$: 6, $\frac{17}{2}$: 2 \\
$(1/2,0)$ & $\langle 3,0\rangle$ & 16 & $\frac{7}{2}$: 1, $\frac{9}{2}$: 2, $\frac{11}{2}$: 5, $\frac{13}{2}$: 5, $\frac{15}{2}$: 2, $\frac{17}{2}$: 1 \\
$(1/2,0)$ & $\langle 1,2\rangle$ & 12 & $\frac{9}{2}$: 2, $\frac{11}{2}$: 4, $\frac{13}{2}$: 4, $\frac{15}{2}$: 2 \\
$(1/2,0)$ & $\langle 3,1\rangle$ & 6 & $\frac{9}{2}$: 1, $\frac{11}{2}$: 2, $\frac{13}{2}$: 2, $\frac{15}{2}$: 1 \\
$(1/2,0)$ & $\langle 1,3\rangle$ & 2 & $\frac{11}{2}$: 1, $\frac{13}{2}$: 1 \\
\midrule
\multicolumn{4}{l}{\textbf{Spin $(1/2,1/2)$}}\\
$(1/2,1/2)$ & $\langle 0,0\rangle$ & 20 & $3$: 1, $4$: 2, $5$: 4, $6$: 6, $7$: 4, $8$: 2, $9$: 1 \\
$(1/2,1/2)$ & $\langle 0,1\rangle$ & 36 & $3$: 1, $4$: 4, $5$: 8, $6$: 10, $7$: 8, $8$: 4, $9$: 1 \\
$(1/2,1/2)$ & $\langle 2,0\rangle$ & 40 & $3$: 1, $4$: 4, $5$: 9, $6$: 12, $7$: 9, $8$: 4, $9$: 1 \\
$(1/2,1/2)$ & $\langle 0,2\rangle$ & 20 & $4$: 2, $5$: 5, $6$: 6, $7$: 5, $8$: 2 \\
$(1/2,1/2)$ & $\langle 2,1\rangle$ & 24 & $4$: 2, $5$: 6, $6$: 8, $7$: 6, $8$: 2 \\
$(1/2,1/2)$ & $\langle 4,0\rangle$ & 4 & $5$: 1, $6$: 2, $7$: 1 \\
$(1/2,1/2)$ & $\langle 0,3\rangle$ & 4 & $5$: 1, $6$: 2, $7$: 1 \\
$(1/2,1/2)$ & $\langle 2,2\rangle$ & 4 & $5$: 1, $6$: 2, $7$: 1 \\
\midrule
\multicolumn{4}{l}{\textbf{Spin $(1/2,1)$}}\\
$(1/2,1)$ & $\langle 1,0\rangle$ & 32 & $\frac{7}{2}$: 2, $\frac{9}{2}$: 5, $\frac{11}{2}$: 9, $\frac{13}{2}$: 9, $\frac{15}{2}$: 5, $\frac{17}{2}$: 2 \\
$(1/2,1)$ & $\langle 1,1\rangle$ & 30 & $\frac{7}{2}$: 1, $\frac{9}{2}$: 5, $\frac{11}{2}$: 9, $\frac{13}{2}$: 9, $\frac{15}{2}$: 5, $\frac{17}{2}$: 1 \\
$(1/2,1)$ & $\langle 3,0\rangle$ & 12 & $\frac{9}{2}$: 2, $\frac{11}{2}$: 4, $\frac{13}{2}$: 4, $\frac{15}{2}$: 2 \\
$(1/2,1)$ & $\langle 1,2\rangle$ & 8 & $\frac{9}{2}$: 1, $\frac{11}{2}$: 3, $\frac{13}{2}$: 3, $\frac{15}{2}$: 1 \\
$(1/2,1)$ & $\langle 3,1\rangle$ & 2 & $\frac{11}{2}$: 1, $\frac{13}{2}$: 1 \\
\midrule
\multicolumn{4}{l}{\textbf{Spin $(1/2,3/2)$}}\\
$(1/2,3/2)$ & $\langle 0,0\rangle$ & 8 & $4$: 1, $5$: 2, $6$: 2, $7$: 2, $8$: 1 \\
$(1/2,3/2)$ & $\langle 0,1\rangle$ & 12 & $4$: 1, $5$: 3, $6$: 4, $7$: 3, $8$: 1 \\
$(1/2,3/2)$ & $\langle 2,0\rangle$ & 12 & $4$: 1, $5$: 3, $6$: 4, $7$: 3, $8$: 1 \\
$(1/2,3/2)$ & $\langle 0,2\rangle$ & 4 & $5$: 1, $6$: 2, $7$: 1 \\
$(1/2,3/2)$ & $\langle 2,1\rangle$ & 4 & $5$: 1, $6$: 2, $7$: 1 \\
\midrule
\multicolumn{4}{l}{\textbf{Spin $(1/2,2)$}}\\
$(1/2,2)$ & $\langle 1,0\rangle$ & 4 & $\frac{9}{2}$: 1, $\frac{11}{2}$: 1, $\frac{13}{2}$: 1, $\frac{15}{2}$: 1 \\
$(1/2,2)$ & $\langle 1,1\rangle$ & 2 & $\frac{11}{2}$: 1, $\frac{13}{2}$: 1 \\
\midrule
\multicolumn{4}{l}{\textbf{Spin $(1,0)$}}\\
$(1,0)$ & $\langle 0,0\rangle$ & 11 & $3$: 1, $5$: 4, $6$: 1, $7$: 4, $9$: 1 \\
$(1,0)$ & $\langle 0,1\rangle$ & 23 & $3$: 1, $4$: 2, $5$: 6, $6$: 5, $7$: 6, $8$: 2, $9$: 1 \\
$(1,0)$ & $\langle 2,0\rangle$ & 20 & $4$: 3, $5$: 3, $6$: 8, $7$: 3, $8$: 3 \\
$(1,0)$ & $\langle 0,2\rangle$ & 8 & $5$: 3, $6$: 2, $7$: 3 \\
$(1,0)$ & $\langle 2,1\rangle$ & 13 & $4$: 1, $5$: 3, $6$: 5, $7$: 3, $8$: 1 \\
$(1,0)$ & $\langle 4,0\rangle$ & 2 & $5$: 1, $7$: 1 \\
$(1,0)$ & $\langle 0,3\rangle$ & 3 & $5$: 1, $6$: 1, $7$: 1 \\
$(1,0)$ & $\langle 2,2\rangle$ & 1 & $6$: 1 \\
\midrule
\multicolumn{4}{l}{\textbf{Spin $(1,1/2)$}}\\
$(1,1/2)$ & $\langle 1,0\rangle$ & 32 & $\frac{7}{2}$: 2, $\frac{9}{2}$: 5, $\frac{11}{2}$: 9, $\frac{13}{2}$: 9, $\frac{15}{2}$: 5, $\frac{17}{2}$: 2 \\
$(1,1/2)$ & $\langle 1,1\rangle$ & 30 & $\frac{7}{2}$: 1, $\frac{9}{2}$: 5, $\frac{11}{2}$: 9, $\frac{13}{2}$: 9, $\frac{15}{2}$: 5, $\frac{17}{2}$: 1 \\
$(1,1/2)$ & $\langle 3,0\rangle$ & 12 & $\frac{9}{2}$: 2, $\frac{11}{2}$: 4, $\frac{13}{2}$: 4, $\frac{15}{2}$: 2 \\
$(1,1/2)$ & $\langle 1,2\rangle$ & 8 & $\frac{9}{2}$: 1, $\frac{11}{2}$: 3, $\frac{13}{2}$: 3, $\frac{15}{2}$: 1 \\
$(1,1/2)$ & $\langle 3,1\rangle$ & 2 & $\frac{11}{2}$: 1, $\frac{13}{2}$: 1 \\
\midrule
\multicolumn{4}{l}{\textbf{Spin $(1,1)$}}\\
$(1,1)$ & $\langle 0,0\rangle$ & 14 & $4$: 2, $5$: 2, $6$: 6, $7$: 2, $8$: 2 \\
$(1,1)$ & $\langle 0,1\rangle$ & 19 & $4$: 2, $5$: 4, $6$: 7, $7$: 4, $8$: 2 \\
$(1,1)$ & $\langle 2,0\rangle$ & 20 & $4$: 1, $5$: 6, $6$: 6, $7$: 6, $8$: 1 \\
$(1,1)$ & $\langle 0,2\rangle$ & 10 & $4$: 1, $5$: 2, $6$: 4, $7$: 2, $8$: 1 \\
$(1,1)$ & $\langle 2,1\rangle$ & 7 & $5$: 2, $6$: 3, $7$: 2 \\
$(1,1)$ & $\langle 4,0\rangle$ & 1 & $6$: 1 \\
\midrule
\multicolumn{4}{l}{\textbf{Spin $(1,3/2)$}}\\
$(1,3/2)$ & $\langle 1,0\rangle$ & 12 & $\frac{9}{2}$: 2, $\frac{11}{2}$: 4, $\frac{13}{2}$: 4, $\frac{15}{2}$: 2 \\
$(1,3/2)$ & $\langle 1,1\rangle$ & 8 & $\frac{9}{2}$: 1, $\frac{11}{2}$: 3, $\frac{13}{2}$: 3, $\frac{15}{2}$: 1 \\
$(1,3/2)$ & $\langle 3,0\rangle$ & 2 & $\frac{11}{2}$: 1, $\frac{13}{2}$: 1 \\
\midrule
\multicolumn{4}{l}{\textbf{Spin $(1,2)$}}\\
$(1,2)$ & $\langle 0,0\rangle$ & 2 & $5$: 1, $7$: 1 \\
$(1,2)$ & $\langle 0,1\rangle$ & 3 & $5$: 1, $6$: 1, $7$: 1 \\
$(1,2)$ & $\langle 2,0\rangle$ & 1 & $6$: 1 \\
\midrule
\multicolumn{4}{l}{\textbf{Spin $(3/2,0)$}}\\
$(3/2,0)$ & $\langle 1,0\rangle$ & 12 & $\frac{7}{2}$: 1, $\frac{9}{2}$: 2, $\frac{11}{2}$: 3, $\frac{13}{2}$: 3, $\frac{15}{2}$: 2, $\frac{17}{2}$: 1 \\
$(3/2,0)$ & $\langle 1,1\rangle$ & 8 & $\frac{9}{2}$: 1, $\frac{11}{2}$: 3, $\frac{13}{2}$: 3, $\frac{15}{2}$: 1 \\
$(3/2,0)$ & $\langle 3,0\rangle$ & 4 & $\frac{9}{2}$: 1, $\frac{11}{2}$: 1, $\frac{13}{2}$: 1, $\frac{15}{2}$: 1 \\
$(3/2,0)$ & $\langle 1,2\rangle$ & 2 & $\frac{11}{2}$: 1, $\frac{13}{2}$: 1 \\
\midrule
\multicolumn{4}{l}{\textbf{Spin $(3/2,1/2)$}}\\
$(3/2,1/2)$ & $\langle 0,0\rangle$ & 8 & $4$: 1, $5$: 2, $6$: 2, $7$: 2, $8$: 1 \\
$(3/2,1/2)$ & $\langle 0,1\rangle$ & 12 & $4$: 1, $5$: 3, $6$: 4, $7$: 3, $8$: 1 \\
$(3/2,1/2)$ & $\langle 2,0\rangle$ & 12 & $4$: 1, $5$: 3, $6$: 4, $7$: 3, $8$: 1 \\
$(3/2,1/2)$ & $\langle 0,2\rangle$ & 4 & $5$: 1, $6$: 2, $7$: 1 \\
$(3/2,1/2)$ & $\langle 2,1\rangle$ & 4 & $5$: 1, $6$: 2, $7$: 1 \\
\midrule
\multicolumn{4}{l}{\textbf{Spin $(3/2,1)$}}\\
$(3/2,1)$ & $\langle 1,0\rangle$ & 12 & $\frac{9}{2}$: 2, $\frac{11}{2}$: 4, $\frac{13}{2}$: 4, $\frac{15}{2}$: 2 \\
$(3/2,1)$ & $\langle 1,1\rangle$ & 8 & $\frac{9}{2}$: 1, $\frac{11}{2}$: 3, $\frac{13}{2}$: 3, $\frac{15}{2}$: 1 \\
$(3/2,1)$ & $\langle 3,0\rangle$ & 2 & $\frac{11}{2}$: 1, $\frac{13}{2}$: 1 \\
\midrule
\multicolumn{4}{l}{\textbf{Spin $(3/2,3/2)$}}\\
$(3/2,3/2)$ & $\langle 0,0\rangle$ & 4 & $5$: 1, $6$: 2, $7$: 1 \\
$(3/2,3/2)$ & $\langle 0,1\rangle$ & 4 & $5$: 1, $6$: 2, $7$: 1 \\
$(3/2,3/2)$ & $\langle 2,0\rangle$ & 4 & $5$: 1, $6$: 2, $7$: 1 \\
\midrule
\multicolumn{4}{l}{\textbf{Spin $(3/2,2)$}}\\
$(3/2,2)$ & $\langle 1,0\rangle$ & 2 & $\frac{11}{2}$: 1, $\frac{13}{2}$: 1 \\
\midrule
\multicolumn{4}{l}{\textbf{Spin $(2,0)$}}\\
$(2,0)$ & $\langle 0,0\rangle$ & 3 & $4$: 1, $6$: 1, $8$: 1 \\
$(2,0)$ & $\langle 0,1\rangle$ & 1 & $6$: 1 \\
$(2,0)$ & $\langle 2,0\rangle$ & 2 & $5$: 1, $7$: 1 \\
$(2,0)$ & $\langle 0,2\rangle$ & 1 & $6$: 1 \\
\midrule
\multicolumn{4}{l}{\textbf{Spin $(2,1/2)$}}\\
$(2,1/2)$ & $\langle 1,0\rangle$ & 4 & $\frac{9}{2}$: 1, $\frac{11}{2}$: 1, $\frac{13}{2}$: 1, $\frac{15}{2}$: 1 \\
$(2,1/2)$ & $\langle 1,1\rangle$ & 2 & $\frac{11}{2}$: 1, $\frac{13}{2}$: 1 \\
\midrule
\multicolumn{4}{l}{\textbf{Spin $(2,1)$}}\\
$(2,1)$ & $\langle 0,0\rangle$ & 2 & $5$: 1, $7$: 1 \\
$(2,1)$ & $\langle 0,1\rangle$ & 3 & $5$: 1, $6$: 1, $7$: 1 \\
$(2,1)$ & $\langle 2,0\rangle$ & 1 & $6$: 1 \\
\midrule
\multicolumn{4}{l}{\textbf{Spin $(2,3/2)$}}\\
$(2,3/2)$ & $\langle 1,0\rangle$ & 2 & $\frac{11}{2}$: 1, $\frac{13}{2}$: 1 \\
\midrule
\multicolumn{4}{l}{\textbf{Spin $(2,2)$}}\\
$(2,2)$ & $\langle 0,0\rangle$ & 1 & $6$: 1 \\
\bottomrule
\end{longtable}

\endgroup

\end{document}
